\section{Agent Benchmarks}

一旦任务需要工具调用、长时推理或多轮迭代，传统单轮 benchmark 就不够了。于是出现了 agent benchmark。

\subsection{SWE-bench}

SWE-bench 要求模型面对真实 codebase 和 issue 描述，最后提交一个 PR，由单元测试来验证是否修复成功。它衡量的是实际软件工程能力，而不只是代码片段生成能力。

\subsection{CyBench}

CyBench 是安全或攻防相关的 CTF 任务，常用 first-solve time 等方式衡量难度和表现。它提醒我们：有些能力一旦做强了，既能用于正当防御，也能用于攻击。

\subsection{MLEBench}

MLEBench 更接近数据科学和机器学习工程任务：清洗数据、训练模型、调参、跑实验。它测的不仅是“会不会答题”，而是“会不会完成一个长流程任务”。

\begin{importantbox}{Agent benchmark 的特点}
Agent benchmark 测的是 \textbf{模型 + 脚手架} 的整体效果。这样更接近真实产品，但也更难分清到底是模型本体改进，还是外部 scaffolding 起了作用。
\end{importantbox}

